% ============================================================
% Section 1: Introduction
% ============================================================

\section{Introduction}
\label{sec:introduction}

Large language models (LLMs) have achieved remarkable capabilities in
text generation, reasoning, and instruction following
\cite{brown2020gpt3, touvron2023llama}. However, these models suffer from
a fundamental limitation: \textbf{they do not know when they do not know}.

The output of a conventional LLM is a logit vector over the vocabulary,
followed by a softmax operation that produces probabilities. While the
entropy of this distribution may serve as a rough proxy for uncertainty,
it conflates aleatoric uncertainty (inherent to the data) with epistemic
uncertainty (due to the model's lack of knowledge). Furthermore, models
trained with cross-entropy tend to be \textbf{poorly calibrated}: the
reported confidence does not correspond to actual accuracy
\cite{guo2017calibration}.

Prior work addresses this problem extrinsically: ensembles
\cite{lakshminarayanan2017simple}, Monte Carlo Dropout
\cite{gal2016dropout}, or orchestration systems that analyze responses
after generation. These approaches add latency, computational costs,
and do not allow the epistemic signal to influence the generation
process itself.

\subsection{Contributions}

In this work, we propose \aletheion{}, a decoder-only LLM that integrates
the epistemic system \textbf{intrinsically} into the neural architecture:

\begin{enumerate}
    \item \textbf{Per-Token Epistemic Tomography}: Each token produces,
    in addition to logits, a set of 30+ epistemic fields organized in
    3 tiers, including decomposed uncertainty ($q_1$, $q_2$), calibrated
    confidence, 5D manifold coordinates, and cognitive state.

    \item \textbf{Differentiable Riemannian Manifold (DRM)}: Each token
    is mapped to a learned 5D space with Riemannian geometry, where
    geodesic distances define epistemic confidence via Metric-Aware
    Distance (MAD).

    \item \textbf{Intentionality Vector (VI)}: A monitoring and
    correction system for manifold health, based on 4 components
    ($\phi$) that detect dimensional collapse and apply deterministic
    corrections.

    \item \textbf{MPC Navigator}: Predictive control on the manifold
    with beam search (K=4, D=3) that plans optimal trajectories over
    12 intervention actions.

    \item \textbf{9 Extension Modules}: EidosDecay, Filosofia3
    ($\phi$-$\psi$ conflict), SelfModel, Grounding, PlasticityGate,
    MPL/Frontier, MOPsi, CausalState, and Metacognitive --- totaling
    {\raise.17ex\hbox{$\scriptstyle\sim$}}364K parameters ($<$0.3\% of the model).

    \item \textbf{Composite Loss with Annealing}: 13 loss components
    with progressive warmup that stabilizes training.

    \item \textbf{Continual Learning}: EWC with online Fisher
    accumulation and Experience Replay with reservoir sampling
    integrated into the trainer.
\end{enumerate}

\subsection{Philosophical Motivation}

The name \textit{Aletheion} derives from the Greek
$\alpha\lambda\acute{\eta}\theta\epsilon\iota\alpha$ (\textit{al\=etheia}),
meaning ``unconcealment'' or ``truth'' in the Heideggerian sense. The
proposal is that the model should not merely generate text, but
\textbf{unconceal} its own epistemic state --- making explicit what it
knows, what it does not know, and how reliable each token is.

\subsection{Paper Organization}

\Cref{sec:architecture} describes the overall transformer architecture.
\Cref{sec:epistemic} details the \texttt{EpistemicHead} and its sub-heads.
Sections~\ref{sec:drm}--\ref{sec:mpc} elaborate on the core components
(DRM, MAD, VI, MPC). \Cref{sec:extensions} covers the 9 extension modules.
\Cref{sec:loss} presents the composite loss. \Cref{sec:continual} discusses
continual learning. \Cref{sec:scalability} addresses scalability.
\Cref{sec:conclusion} concludes.
